\begin{abstract}
    Text-to-brick generation requires translating natural-language descriptions into spatial arrangements of discrete pieces. Concise prompts leave object dimensions, proportions, and part layouts implicit, making geometric data construction and structural constraints central to this task. We present BrickForma, a framework combining voxel-to-brick data construction, language-model fine-tuning, and geometric reconstruction. LLM-assisted parametric programs generate 80,000 source voxel models across 21 categories on a $48^3$ grid, each paired with five levels of description. The resulting approximately 400,000 caption records link varying amounts of linguistic detail to shared object geometry. A conversion pipeline integrates same-color tiling, connection analysis, and local bridging and retiling to construct geometrically checked brick targets. We fine-tune Qwen3-1.7B with low-rank adaptation on caption--placement pairs to predict brick dimensions and positions from text. A subsequent reconstruction stage uses mixed-integer linear programming to seek an exact cover of the reconstructed voxel target with a single stud-connected component. Experiments compare BrickForma with BrickGPT and BrickNet on 114 shared prompts spanning five description levels. Under a common rendering protocol, BrickForma achieves the highest mean VQAScore, PE, and SigLIP2 at every level, including the shortest descriptions, while matching BrickGPT's mean invalid-placement ratio of 0.00\%.
\end{abstract}
